\documentclass[12pt]{article}
\usepackage[margin=1.0in]{geometry}
\usepackage{setspace}
\setstretch{2}
\usepackage{graphicx}

\usepackage{hyperref}
\hypersetup{%
    colorlinks=true,
    linkcolor=cyan,
    filecolor=magenta,
    urlcolor=blue,
    citecolor=blue,
}
\usepackage{siunitx}
\DeclareSIUnit{\molar}{M} 

% Bibliography
\usepackage{natbib}
\bibliographystyle{msb3} % the optional note field in the bib file is useful for MSB data refs

\begin{document}

\begin{titlepage}
    \begin{center}
        \vspace*{1cm}

        {\Large Atomic elementary flux modes explain the steady state flow of
        metabolites in large-scale flux networks}

%        \vspace{1.0cm}
%        aEFMs explain metabolite flows
            
        \vspace{1.0cm}

        Justin G. Chitpin$^{1,2}$ and Theodore J. Perkins$^{1,2*}$

        \vspace{1.0cm}

        $^1$Ottawa Hospital Research Institute, 501 Smyth Road, Ottawa, K1H 8L6, Ontario, Canada\\
        $^2$Ottawa Institute of Systems Biology, Department of Biochemistry, Microbiology and Immunology, University of Ottawa, 451 Smyth Road, Ottawa, K1H 8M5, Ontario, Canada
           
        \vspace{0.5cm}

        $^*$ Corresponding author.

        \vspace{1.0cm}
     
        Flux network/Elementary flux mode/flow decomposition/Markov chain/atom mapping
            
   \end{center}
\end{titlepage}

\clearpage

\section*{Abstract} % 175 words max

Steady state fluxes are a measure of cellular activity under metabolic
homoeostasis, but understanding how individual substrates are metabolized
remains a challenge in large-scale networks. Pathway-based approaches such
as elementary flux mode (EFM) analysis are limited to small networks due
to the combinatorial explosion of pathways and the ambiguity of
decomposing fluxes onto EFMs. Here, we present an alternative approach to
explain metabolic fluxes in terms of the steady state flow of their atomic
constituents. We refer to these pathways as atomic elementary flux modes
(AEFMs) and show that computations involving AEFMs are orders of magnitude
faster than standard EFMs. Using our approach, we enumerate carbon and
nitrogen AEFMs in five genome-scale metabolic models and compute the AEFM
decomposition of fluxes estimated in a HepG2 liver cancer cell line. Our
results systematically characterize carbon and nitrogen remodelling and,
on the HepG2 network, predict glutamine metabolism through a recently
discovered non-canonical TCA cycle.

\clearpage

\section*{Introduction} % intros tend to be very succinct

% (a) flux overview
Recent technological and computational advancements are leading to a rise in
metabolic flux
data~\citep{SternAlon2023Imac,LeeWonDong2019Spas,WangLin2022Srit,WangRuohong2022Gstm}.
These fluxes quantify the rate of metabolite interconversions within
a cell or their movements across subcellular compartments. Metabolic
fluxes are studied to understand the dynamics of substrate
switching~\citep{ZampieriMattia2012Rmuc,KochanowskiKarl2021Gcom}, redox
balancing~\citep{McKinlayJamesB2010Cdfa,FranchinaDavideG2022Grbc},
overflow
metabolism~\citep{BasanMarkus2015OmiE,VanderHeidenMatthewG2009UtWE}, cell
differentiation~\citep{EndoYusuke2022Fami,BenjaminJacksonT2024Mrot}, and
cellular
communities~\citep{SchuetzRobert2012MOoM,YuJasonSL2022Mcfr,Gustafsson2024Mcbc}.
Hence, `fluxomics'~\citep{EmwasAbdul-Hamid2022F-NM,SanfordKarl2002Gtfa} is
emerging as a field dedicated to quantifying the totality of fluxes within
biological systems.

% (b) estimating steady state fluxes - happy with the writing
Current experimental and computational methods generally estimate fluxes
under metabolic homeostasis. Under this assumption, the production and
consumption of metabolites are balanced over time, and the fluxes are said
to be under steady state. Techniques such stable isotope tracing analysis
(SITA) experimentally determine steady state fluxes by culturing cells in
labelled substrates and quantifying their incorporation within endogenous
metabolites~\citep{BuescherJoergM2015Arfi,AntoniewiczMaciekR2018Agt1}.
Using this method, fluxes have been measured in major metabolic subsystems
such as glycolysis, pentose phosphate pathway, and the TCA
cycle~\citep{OrthJeffreyD2010RaUo,AhnEunyong2017Tfro,WuChao2019FARC}. Many
computational methods, including but not limited to flux balance analysis
(FBA), have also been proposed to infer steady state fluxes from transcript,
protein, and metabolite abundance
data~\citep{UematsuSaori2022Mlmf,GonzalezArrueNicolas2023Peom,HuangYuefan2023Ccmf,KasteJoshuaAM2023Afpu,LeeJustinY2023Tiac}.
Collectively, these approaches are being used to estimate and predict fluxes in
genome-scale metabolic models (GEMs) involving hundreds to thousands of
reactions~\citep{NilssonAvlant2020Qaoa,CherkaouiSarah2022Afao,GelbachPatrickE2024Fsig}.

% (c) introducing problem
Although much effort has been made to estimate steady state fluxes, less
emphasis has been placed on developing methods to analyze these
large-scale flux datasets. In contrast to abundance data (e.g.\
transcripts, proteins, metabolites), steady state fluxes are constrained
by the reaction stoichiometries between metabolite inflows and outflows.
In a network with fixed source and sink fluxes, an increase in flux along
one metabolic pathway must decrease fluxes along one or more other
pathways to maintain mass conservation. For simple networks, one may
predict how the system would evolve by identifying the up and downstream
reactions flanking an enzymatic perturbation. For example,
\citeauthor{SchwartzJean-Marc2006Qema} were able to
enumerate all glycolytic pathway fluxes in a 12-reaction model of yeast
glycolysis~\citep{SchwartzJean-Marc2006Qema}. Yet in more complex genome-scale
metabolic networks containing thousands of fluxes
(e.g.\ \cite{BrunkElizabeth2018Reat,RobinsonJonathanL2020Aaoh,WangHao2021Gmnr}),
there exists no hierarchical relationship between reactions. A change in one
reaction rate, let alone multiple perturbations, may alter network fluxes
beyond neighbouring reactions~\citep{NobileMarcoS2021Agsa}.

% (d) EFMs
This problem of understanding the flow of metabolites within steady state
flux networks has historically been addressed through the concept of
elementary flux modes (EFMs). EFMs are a pathway-based approach to analyze
metabolic networks and are defined as minimal sets of biochemical
reactions that maintain steady state flux~\citep{SchusterStefan1994Oefm}.
This minimal property specifies that an EFM cannot be decomposed into two
or more smaller pathways carrying steady state flux. An attractive
property of EFMs is that any set of steady state fluxes in a network can
be explained as a positive, linear combination of
EFMs~\citep{SchusterStefan1999Doef}. While EFMs are a mechanistic
explanation of how metabolites travel through networks, the number of EFMs
scales combinatorially as a function of network
size~\citep{AcunaVicente2009Maci}. It remains computationally infeasible
to enumerate EFMs in GEMs after over a decade of algorithmic advancements
(e.g.\
\citealt{GagneurJulien2004Coem,TerzerMarco2008Lcoe,BuchnerBiancaA2021EaPp}).
The number of EFMs also poses a challenge for explaining network fluxes in
terms of their EFMs. When there are more EFMs than linearly independent
fluxes, there exist multiple ways to assign EFM weights that reconstruct
the observed network fluxes. Consequently, structural analyses of EFMs and
downstream analyses of their weights have been limited to small-scale
metabolic
subnetworks~\citep{FakihIbrahim2023Dgmm,MpabanziLiliane2022Frca,ToyaYoshihiro2022Mpef,MatsudaFumio2011Esoy,SchwartzJean-Marc2006Qema}.

In our previous work, we developed a novel solution to the EFM flux
decomposition problem for the special case of unimolecular reaction
networks under a Markov constraint. We modelled a hypothetical particle
moving between metabolites under the assumption that the transition
probability of each reaction was proportional to the observed steady state
flux. By expanding the notion of Markov chain state to include sequences
of metabolites, we proposed an algorithm to efficiently compute simple
cycle probabilities corresponding to a given EFM. We called this construct
the \emph{cycle-history Markov chain} (CHMC) and proved that these EFM
probabilities could be scaled to weights that fully reconstructed the flux
along each reaction~\citep{ChitpinJustinG2023AMct}.

% (e) directives
In this work, we generalize our CHMC method to operate on any type of
metabolic network, including those involving multispecies reactions. We do
so by proposing a new type of \emph{atomic elementary flux mode} which
traces the steady state flow of an individual atom within a metabolic
network. We refer to these pathways as AEFMs to distinguish them from the
(molecular) EFMs proposed by \citeauthor{SchusterStefan1994Oefm}. While
a linear combination of molecular EFMs will reconstruct the overall
reaction fluxes, the totality of AEFM weights for a given source
metabolite will always reconstruct its mass flow entering the network. We
first introduce our framework for enumerating and uniquely identifying
AEFM weights. Using a state-of-the-art atom mapping algorithm, we show how
to generate atomic transition graphs from which we construct atomic
cycle-history Markov chains (ACHMCs). From these ACHMCs, we enumerate
AEFMs in GEMs considered infeasible by state-of-the-art molecular EFM
enumeration programs. We further analyse the structure of the AEFMs to
characterize nutrient source remodelling. Lastly, we analyze a single set
of AEFM weights estimated from a HepG2 liver cancer cell line to quantify
glutamine metabolism. Our results predict glutamine in HepG2 cells is
metabolized through a non-canonical TCA pathway recently discovered in
non-small cell lung cancer, mouse embryonic, and C2C12 mouse myoblast
cells~\citep{ArnoldPaigeK2022Anta}.

\section*{Results}

\subsection*{Tracing atomic flows in genome-scale models} % pipeline overview

We first describe what molecular EFMs look like in metabolic networks and
provide intuition why enumerating and decomposing fluxes onto these
pathways remains a fundamental problem. Figure \ref{fig1Pipe}a shows an
example molecular EFM involving substrates glucose and ATP which are
consumed to produce metabolites G3P, DHAP, and ADP. Multiple source and
sink fluxes are required to balance the internal metabolite
stoichiometries, leading to several branching and converging paths within
the molecular EFM. As the metabolic networks expands to include dozens or
more reactions, the corresponding molecular EFMs can grow increasingly
complex with numerous source and sink fluxes required to
stoichiometrically balance participating metabolites. In particular, there
may be multiple sets of reactions involved to balance pairs of energetic
molecules (e.g.\ ATP/ADP), redox equivalents (e.g.\ NAD$^+$/NADH), metabolic
byproducts (e.g.\ water, hydrogen ions), or coenzymes (e.g.\ coenzyme A).
This complexity can lead to an intractable number of pathways that may
only differ by reaction subsets required to mass-balance these
metabolites. Choosing certain metabolites to remove may, in part, solve
this problem. However, the choice of metabolites to prune from the network
is often subjective and not guaranteed to sufficiently reduce the total
number of molecular EFMs.

Given many metabolic flux networks are generated to study nutrient
catabolism~\citep{ZamboniNicola2015DtMS} or engineer desirable metabolic
products~\citep{ZanghelliniJurgen2013Efmi}, we argue for studying source
metabolite-derived AEFMs to reduce the combinatorial number of molecular
EFMs. We highlight an example carbon AEFM in Figure \ref{fig1Pipe}b which
shows how the C6 glucose carbon is remodelled during the first steps of
glycolysis. As atoms are never split apart or fused together in metabolic
reactions, these AEFMs always correspond to simple paths spanning source
and sink metabolites or simple cycles between internal metabolites. These
pathways are straightforward to interpret because each atom within an AEFM
can only transit from a single substrate to product, regardless of the
network size, connectivity, or number of participating substrates/products
in each reaction. This property further enables us to enumerate and
compute the AEFM weights using our previously described CHMC method to
uniquely decompose steady state fluxes onto EFMs for unimolecular reaction
networks~\citep{ChitpinJustinG2023AMct}.

We outline our computational pipeline in Figure \ref{fig1Pipe}c to
enumerate AEFMs and uniquely assign their weights when the steady state
network fluxes are known. In step 1, we require as inputs the metabolic
network defined by a set of metabolites and reactions encoded as
a stoichiometry matrix, the metabolite structures denoted as SMILES
strings, and the steady state fluxes. From the stoichiometry matrix and
SMILES strings, we construct reaction SMILES strings which are imputed
into the atom mapping program RXNMapper~\citep{Schwaller2021extraction} to
map the position of each atom within each reaction equation. This atom
mapping data is used to construct an atomic transition graph in step 2,
which traces the atomic position of an individual atom from a given source
metabolite entering the network. We set the transition probabilities
proportional to the atomic flux through each reaction, accounting for
those with multiple substrate stoichiometries (see Methods and Figure
\ref{figSaccr}).
%For example, 10\SI{}{\milli\molar} of input glucose flux would correspond
%with 10\SI{}{\milli\molar} of atomic flux through each glucose carbon.
For each atomic transition graph, we construct the corresponding ACHMC to
enumerate all AEFMs involving that source metabolite atom in step 3.
Finally, the AEFM weights are computed in step 4 by steady state analysis
of the ACHMC. 

\subsection*{Enumerating AEFMs is computationally tractable in large-scale networks}

Given the major computational challenge of enumerating molecular EFMs in
networks with $>100$ metabolites and reactions, we first sought to test
whether it was feasible to enumerate AEFMs in GEMs. We selected five
genome-scale metabolic models across distinct organisms in ascending order
of network size. These networks include a metabolic model of \textit{E.
coli} (E. coli core; \citealt{OrthJeffreyD2010RaUo}; Data ref:
\citealt{OrthJeffreyD2010RaUoDATASET}), human red blood cells (iAB RBC
283; \citealt{BordbarAarash2011iApd}; Data ref:
\citealt{BordbarAarash2011iApdDATASET}), \textit{H. pylori} (iIT341;
\citealt{ThieleInes2005EMRo}; Data ref:
\citealt{ThieleInes2005EMRoDATASET}), \textit{S. aureus} (iSB619;
\citealt{BeckerScottA2005Grot}; Data ref:
\citealt{BeckerScottA2005GrotDATASET}), and human liver cancer cells
(HepG2; \citealt{NilssonAvlant2020Qaoa}; Data ref:
\citealt{NilssonAvlant2020QaoaDATASET}). These models were pre-processed
to ensure unambiguous atom mappings across all reactions. Namely, all
metabolites with no known structures (pseudometabolites), and
pseudoreactions with non-integer stoichiometries were removed from the
networks. The resulting metabolic networks ranged between 71-547
metabolites and 144-974 reactions (Figure \ref{figSgh}). We then attempted to
enumerate the molecular EFMs using FluxModeCalculator, which is the
fastest molecular EFM enumeration
program~\citep{vanKlinkenJanBert2016Faet}. Following our pipeline, we
subsequently compiled the metabolite SMILES strings and constructed ACHMCs
rooted on each carbon and nitrogen atom across all source metabolites. The
resulting summary statistics describing the GEMs and ACHMCs are presented
in Table 1. Across all carbon ACHMCs, we observed a broad increase in the
total number of AEFMs as a function of metabolic network size. The one
exception to this observation was the HepG2 network which contained fewer
metabolites and reactions than iSB619, yet exhibited two orders of
magnitude more carbon AEFMs. Inspection of the HepG2 network revealed it
only contained 40 source metabolites compared to the 198 source
metabolites in the iSB619 network, indicating a greater degree of network
connectivity between internal metabolites (Figure \ref{figSgh}).

Figure \ref{fig2Bench}a shows the total number of atomic and molecular
EFMs enumerated by our program MarkovWeightedEFMs.jl and
FluxModeCalculator on the five GEMs. Using our ACHMC pipeline, we
enumerated between $10^3 - 10^6$ carbon and $10^2-10^3$ nitrogen AEFMs
within each GEM. In contrast, FluxModeCalculator failed to enumerate
molecular EFMs in all but the two smallest E. coli core and iAB RBC 283
datasets, returning $10^3 - 10^4$ times more molecular-versus-atomic EFMs. 

By focusing on an atom of interest, our ACHMC approach greatly reduced the
number of EFMs to biologically relevant pathways involving nutrient source
remodelling. This reduction in pathways resulted in AEFM-versus-EFM
enumeration completing $10^3-10^4$ times faster (Figure \ref{fig2Bench}b).
While running FluxModeCalculator in parallel did improve run times by
8.2-8.6 times (Figure \ref{figSbfp}), it remained over $10^1-10^2$ slower
  than enumerating their corresponding carbon and nitrogen AEFMs. We
  further found that AEFM enumeration actually involved fewer
  metabolite-atom combinations than the total number of metabolites in
  a given network. Constrained by the atom mapping predictions, we
  observed that the average carbon and nitrogen ACHMC across all five GEMs
  only transited 0.41\% and 1.7\% of the total number of metabolite-carbon
  and metabolite-nitrogen positions, respectively (Figure \ref{figSass}).
  While serial run times in Figure \ref{fig2Bench}c scaled quadratically
  (Figure \ref{fig2Bench}c; $\text{slope}=1.90, \text{intercept}=-6.54$),
  our program and AEFM approach benefits greatly from parallelism. Each
  ACHMC can be computed independently of one another and our program is
  also parallized on the level of individual ACHMCs to improve AEFM
  enumeration. For the top 28 largest carbon ACHMCs in the HepG2 network,
  Figure \ref{fig2Bench}d shows a median speedup of 7.1 times when using
  32 threads. 


\iffalse

These fast AEFM enumeration times may seem counterintuitive given the
number of metabolite-atom positions is much larger than the total number
of metabolites in a given metabolic network. To resolve these findings, we
enumerated the total number of metabolite-carbon and metabolite-nitrogen
positions across all five GEMs. The HepG2 and iAB RBC 283 networks
exhibited the greatest number of combinations with 8194 metabolite-carbon
and 2533 metabolite-nitrogen positions. This corresponded to over 18 and
4 times the number of metabolites within each respective network. However,
  each atomic transition graph may only involve a subset of the entire
  list of metabolite-carbon or metabolite-nitrogen positions. Constrained
  by the atom mapping predictions, we found that the average carbon and
  nitrogen ACHMC across all five GEMs only transited 0.41\% and
1.7\% of the total number of metabolite-carbon and metabolite-nitrogen
states, respectively (Figure \ref{figSass}). Hence, the metabolite-atom state space
of any given ACHMC is several orders of magnitude smaller than the total
number of network metabolites.

We therefore hypothesized that the speedup between MarkovWeightedEFMs.jl
and FluxModeCalculator was due to the proportionally fewer number of
AEFMs-versus-molecular EFMs across the metabolic networks. For molecular
EFM enumeration algorithms, it is well established that run times cannot
be predicted from the number of metabolites or reactions within the
model~\citep{UllahEhsan2020Tsef}. We found, however, that AEFM run times
scale approximately quadratically as a function of the top 28 largest
ACHMCs (Figure \ref{fig2Bench}c; $\text{slope}=1.90, \text{intercept}=-6.54$). While
future work is required to improve the algorithm scaling, we emphasize
that the computation of each ACHMC is embarrassingly parallel at the
source metabolite-atom level. Our software is also parallelizable on the
level of individual ACHMCs to further speed up AEFM enumeration. For those
top 28 largest carbon ACHMCs in the HepG2 network, we re-enumerated their
AEFMs with varying numbers of threads. Figure \ref{fig2Bench}d reports the algorithm
scaling, where we observed a median speedup of 7.1 times when using 32
threads. As expected, the greatest speedups were associated with the
largest ACHMCs. For example, the largest ACHMC corresponded to a glutamine
carbon with 1,120,290 ACHMC states/168,793 AEFMs, which exhibited a 9.3
times speedup when enumerated with 32 threads.

\fi

\subsection*{Most AEFMs are source-to-sink pathways spanning dozens of reactions}

We next turned to characterizing the structural properties of all carbon
and nitrogen AEFM across the five GEMs. Molecular EFMs can be divided into
pathway that span source and sink metabolites or internal loops within the
metabolic network such as pairs of reversible reactions. Both pathways
have clear biological interpretations with source-to-sink pathways
explaining the steady state flow of source metabolites on their way to
exiting the network, and internally looped pathways corresponding to
substrate cycles. At an atomic level, these pathways can be subclassified
further into AEFMs that transit sets of distinct metabolites or those that
revisit the same metabolite, albeit in different atomic positions. This notion
of a ``metabolite revisitation'' is similar to an isotopomer in stable isotope
tracing~\citep{AntoniewiczMaciekR2018Agt1}. An example source-to-sink AEFM
with this property is a glucose-derived carbon AEFM remodelling through
two turns of the TCA cycle; on the second turn, the position of the carbon
atom shifts within the TCA intermediates, resulting in different
isotopomers~\citep{DuanLikun20221tas} and therefore a source-to-sink
pathway with a metabolite revisitation under our nomenclature.

To obtain a broad understanding of the types of AEFMs within the five
GEMs, we categorized them into source-to-sink pathways and looped pathways
with or without metabolite revisitations in different atomic positions.
Figure \ref{fig3Struct}a-b shows the counts for each AEFM classes within
each GEM. As expected, we observed that the majority of carbon AEFMs were
source-to-sink pathways with or without metabolite revisitations. Those
without metabolite revisitations tended to scale as a function of network
size, while those with revisitations increased modestly. Investigation of
these source-to-sink pathways confirmed that many metabolite revisitations
involved TCA cycle intermediates. Looped AEFMs without metabolite
revisitations primarily corresponded to transport reactions cycling
metabolites across compartments and cyclic pathways involving TCA
intermediates. These two observations explain, in part, why the iAB RBC
283 network contained the fewest number of pathways with metabolite
revisitations, since they lack organelles and cannot support cyclic AEFMs
associated with metabolite transport between compartments nor
mitochondrial TCA activity. There were very few internally looped AEFMs
with metabolite revisitations with no obvious trend across the five GEMs.
The greatest number of these pathways were found in the E. coli core
network with many metabolite revisitations involving carbon dioxide
derived from isocitrate and alpha-ketoglutarate re-entering the TCA cycle
via oxaloacetate. As there are fewer metabolites containing nitrogen, we
found the majority of their AEFMs corresponded to pathways without
metabolite revisitations, primarily involving amino acid remodelling or
transports between compartments. Looped pathways with metabolite
revisitations were only observed in the iIT341 and iSB619 networks through
reactions involving the production and consumption of ammonia. 

Since many looped AEFMs involve reversible reactions or metabolite
transport across compartments, we hypothesized that looped AEFMs should
involve fewer reactions than source-to-sink pathways which can span
numerous internal metabolites. We calculated the fraction of carbon and
nitrogen AEFMs with given lengths across all five GEMs and plotted their
distributions in Figure \ref{fig3Struct}c-d. Our results show that carbon
AEFMs in the first four networks are bimodally distributed with the
majority of pathways involving a dozen or fewer reactions. While carbon
AEFMs of the HepG2 network appeared to exhibit a single Gaussian
distribution ($\mu=61.7$, $\sigma^2=15.5$), closer inspection of the
histogram revealed a small peak spanning 2-12 reactions. To determine
whether these peaks were associated with source-to-sink and looped AEFM
classes, we reclassified the carbon AEFMs spanning less than 5-12
reactions depending on the GEM (shorter pathways), and the remaining
carbon AEFMs (longer pathways). We found that longer pathways were
associated with source-to-sink pathways, with shorter pathways containing
the majority of looped pathways. In fact, 65-95\% of all longer pathways
in each GEM were source-to-sink pathways with metabolite revisitations,
with this statistic increasing to 95-100\% when including source-to-sink
pathways without metabolite revisitations. Among shorter pathways, we
observed that roughly half or more of them were looped carbon AEFMs with
or without metabolite revisitations. Analysis of the nitrogen AEFMs
revealed only a single distribution of relatively short pathways, since
there are much fewer nitrogen-containing metabolites across all GEMs.

We next investigated how atoms within the same source metabolite may
follow different paths through the network, irrespective of their AEFM
class. We first computed the number of AEFMs for each carbon atom (resp.
nitrogen atom). For each metabolite, we then computed the ratio of the
class. We then computed the number of AEFMs for each carbon atom (resp.
nitrogen atom). So for example, a ratio of 10 means that one of the carbon
atoms in the metabolite can transit the network in 10 times as many
different ways as the carbon with the fewest paths through the network.
Our results are shown in Figure \ref{fig3Struct}e-f, where each dot
corresponds to the ratio for a given source metabolite. We found that
a surprising number of carbons and nitrogens can transit different
pathways within their respective source metabolites. Focusing on the
source metabolites with carbon and nitrogen ratios over 100, we observed
many of these corresponded to hydrophobic and positively-charged amino
acids such as phenlylalanine, methionine, arginine, and histidine. Source
metabolites containing fewer carbons or nitrogens tended to exhibit ratios
closer to
1 with AEFMs involving the same metabolite subsets. Metabolites such as
  glycine, for example, are equivalently metabolized across the same
  carbon AEFMs. For each GEM, and not counting source metabolites with
  ratios of one, we computed median carbon ratios between
  2.0 and 46.8 and nitrogen ratios between 1.3 and 11.2. Overall, these
  findings highlight the varying degrees of metabolic flexibility for
  atomic constituents of source metabolites to transit the metabolic
  network.

\subsection*{A minority of carbon AEFMs explains the majority of source metabolite remodelling}

% Glutamine carbon 1:   8,098
% Glutamine carbon 2: 168,793
% Glutamine carbon 3:  42,691
% Glutamine carbon 4:  16,581
% Glutamine carbon 5:   4,922
% Glutamine carbon AEFM total: 241,085
% Glutamine carbons can remodel into 202 distinct metabolites aside from
% input glutamine

We finally highlight the potential of our AEFM method for quantifying
nutrient source metabolism. We turned to our HepG2 network which contained
a single set of steady state fluxes estimated by
\citeauthor{NilssonAvlant2020Qaoa}. These fluxes were estimated from
a parsimonious FBA model of HepG2 cells cultured under exponential growth
with external fluxes constrained by time-resolved exometabolomic
measurements (see Methods for more information). Our focus was to analyze
this flux network using our AEFM framework to identify alternative
pathways of glutamine metabolism since many cancer cells, including HepG2,
are dependent on this non-essential amino acid for survival and
proliferation~\citep{YeYanyan2023Gmri,YooHeeChan2020Gric}. Using our
method, we decomposed the set of steady state fluxes onto the 241,085
glutamine carbon AEFMs characterizing the HepG2 metabolic network. Figure
\ref{fig4Hepg2}a shows the cumulative explained mass flow of each
glutamine carbon as a function of their AEFMs ranked from greatest to
smallest explained mass flow. We capped the cumulative explained mass flow
of each glutamine carbon at 99\% because we observed that 98.9\% of AEFMs
collectively explained the remaining 1\% glutamine carbon flow. These
results suggested that the majority of glutamine carbon remodelling is
explained by a surprisingly small number of glutamine AEFMs, especially
since our Markovian flux decomposition approach assigns some weight to all
network pathways. The top 5 ranked AEFMs across each glutamine carbon (25
pathways altogether) explained over half (50.2\%) of the total glutamine
carbon mass flow with the subsequent 5 next highest ranked AEFMs only
explaining an additional 9.8\% glutamine carbon mass flow. These
observations were consistent across carbon AEFMs derived from other amino
acids as well as glucose, leading us to hypothesize that steady state
pathway fluxes are centralized to a relatively small subset of AEFMs
(Figure \ref{figSccl}). For some source metabolites, we observed
overlapping curves between carbons (e.g.\ arginine, glycine, serine in
Figure \ref{figSccl}), suggesting that individual carbons were remodelled
through the same AEFMs. However, corroborating results from Figure
\ref{fig3Struct}e-f, we observed many non-overlapping curves for the
majority of amino acids since their carbons could transit the network via
different internal metabolites.

%We denote this carbon mass flow as $C_k = w_k \cdot L_k$, where $L_k$ is
%the length of EFM $k$. 

To obtain a high-level understanding of the dominant pathways of glutamine
carbon remodelling, we constructed a metabolic subnetwork from the top
5 glutamine carbon AEFMs, collapsing metabolite-atom positions onto
distinct metabolites, since none of these AEFMs involved metabolite
revisitations. Of the 435 metabolites in the HepG2 dataset, input
glutamine can remodel into 202 species localized to the cytosol or
mitochondria. We found that the top 5 glutamine-derived carbon AEFMs
consist of just 23 metabolites and 31 reactions (Figure \ref{fig4Hepg2}b).
These AEFMs generally consisted of cyclic pathways involving TCA
intermediates. Only two source-to-sink reactions convert cytosolic
glutamine towards the protein pool and carbon dioxide through
decarboxylation of alpha-ketoglutarate in the mitochondria.

As AEFMs are steady state pathways, they can be visualized as Sankey
diagrams to better understand the carbon flows through the network. Figure
\ref{fig4Hepg2}c-e shows the top 5 AEFMs for each glutamine carbon. Each
AEFM in the diagram is shown as a flow proportional to its AEFM weight
with nodes representing the metabolite transited by the given glutamine
carbon. Although each glutamine carbon can transit a different number of
AEFMs, we found the top 5 backbone carbons and first two sidechain carbons
were identical in both AEFM pathway and weight, with no AEFM containing
a metabolite revisitation.

Focusing on either backbone carbon in Figure \ref{fig4Hepg2}c, we observed
that all AEFMs corresponded to well-established metabolic subsystems. The
first, third, and fifth highest ranked AEFMs corresponded to metabolites
and reactions involved in the malate-aspartate shuttle. The fourth largest
AEFM resembled a combination of canonical TCA cycle activity and
malate-aspartate shuttle, with TCA intermediates citrate, isocitrate, and
alpha-ketoglutarate localized to the cytosol rather than mitochondria.
This AEFM also involved additional reactions transforming
alpha-ketoglutarate into glutamate before remodelling into succinyl-CoA in
the mitochondria. We noticed that the second largest AEFM cycled carbon
mass between TCA intermediates citrate, oxaloacetate, and malate between
the cytosol and mitochondria. Interestingly, this metabolic pathway
corresponds to the recently discovered non-canonical TCA activity
described by \citeauthor{ArnoldPaigeK2022Anta}, who validated this pathway
in non-small cell lung cancer, mouse embryonic, and C2C12 mouse myoblast
cells by stable isotope tracing analysis~\citep{ArnoldPaigeK2022Anta}.
This non-canonical TCA activity was further observed in glutamine carbons
3 and 4 in AEFM 2 (Figure \ref{fig4Hepg2}d). We also found AEFMs 1 and
3 in Figure \ref{fig4Hepg2}d-e corresponded to the malate-asparate
shuttle. Only the terminal glutamine carbon in Figure \ref{fig4Hepg2}e
contained source-to-sink AEFMs that decarboxylated glutamine-derived
alpha-ketoglutarate into carbon dioxide.

Our analyses revealed that the most active glutamine-derived carbon AEFM
weights corresponded to well-known metabolic subsystems. This finding led
us to ask whether we could summarize our AEFM analysis as a simplified
metabolic model describing the majority of glutamine carbon remodelling in
the HepG2 flux network. The resulting schematic is shown in Figure
\ref{fig4Hepg2}f which consists of the following three major metabolic
pathways: canonical TCA-like cycle, non-canonical TCA cycle, and the
malate-aspartate shuttle. These metabolic pathways were chosen based on
the number of top glutamine-derived carbon AEFMs that overlapped with
reactions within these subsystems. Our results suggest input glutamine is
converted into cytosolic glutamate which is transported into the
mitochondria through aspartate/glutamate antiporters encoded by
\textit{SLC25A12} and \textit{SLC25A13} annotated in the HepG2 metabolic
model. Mitochondrial glutamate is then converted into TCA intermediate
alpha-ketoglutarate and citrate. Instead of following the canonical
reactions of the TCA cycle, the mitochondrial citrate is exported to the
cytosol where it can participate in non-canonical TCA cycle activity or
the malate-aspartate shuttle to re-enter the mitochondria. Taken together,
the dominant carbon AEFMs predicted by our model suggest HepG2 glutamine
is conserved within TCA intermediates and regenerated through subsequent
reactions involving cytosolic citrate.

\section*{Discussion}

Since the proposal of molecular EFMs, much effort has been made to
enumerate and decompose steady state metabolic fluxes onto these pathways.
By defining EFMs with respect to individual atoms rather than molecules,
we introduced the concept of AEFMs which describe how individual atoms are
remodelled within a metabolic network. As AEFMs never involve branching
pathways, they are easier to interpret than their molecular EFM
counterparts, which may involve multiple inputs and outputs to balance
molecular stoichiometries. Using the atom mapping program RXNMapper, we
adapted our previously described CHMC method to efficiently enumerate
carbon and nitrogen AEFMs in five GEMs. Our structural analysis
systematically characterized source metabolite remodelling at the atomic
level. By comparing the number of AEFMs across atoms within the same
source metabolite, we revealed how certain amino acids displayed
remarkable variability in the number of pathways transited by their
constituent atoms. This analysis, in particular, may be useful for
designing non-uniformly labelled stable isotope tracers to identify fluxes
within specific metabolic subsystems. Our AEFM weight analysis of
the HepG2 network revealed that the majority of glutamine carbons were
remodelled via cyclic pathways regenerating TCA intermediates through
cytosolic citrate. Excitingly, our model also predicted non-canonical TCA
activity which has not been previously described in HepG2 cells.

As a systematic method to explain the flow of metabolites, the problem of
molecular EFM enumeration has historically limited their applicability to
small-scale networks. Many algorithmic advances have incrementally
improved molecular EFM enumeration, yet never feasibly beyond networks
with $>100$ metabolites and
reactions~\citep{KampAxelvon2006M5fa,TerzerMarco2008Lcoe,vanKlinkenJanBert2016Faet}.
Molecular EFM enumeration is generally regarded as infeasible for
large-scale networks---so much so that numerous methods have been
developed to enumerate a subset of molecular EFMs constrained by pathway
length~\citep{deFigueiredoLuisF2009Ctse}, reaction
thermodynamics~\citep{GerstlMatthiasP2015Mief}, participating
metabolites~\citep{PeyJon2014Dcoe}, participating
reactions~\citep{KaletaChristoph2009Ctwb}, gene regulatory
rules~\citep{JungreuthmayerChristian2013rSue}, metabolic
subsystems~\citep{SchusterS2002Etps,SchwartzJean-Marc2007Omfa},
extracellular flux
measurements~\citep{SoonsZitaITA2011Iomm,JungersRaphaelM.2011Fcom}, or
even random sampling~\citep{MachadoDaniel2012Rsoe}. While our HepG2
results did show the majority of source metabolite fluxes were explained
by a small number of AEFMs, contextualizing the significance of these
dominant pathway fluxes is only possible when all AEFMs are known.
Furthermore, reaction thermodynamics, gene regulatory rules, and other
constraints may change across different biological conditions. By
computing AEFMs, instead, we could feasibly enumerate atomic pathways in
genome-scale networks containing hundreds of metabolites and reactions
within minutes or hours on consumer computing hardware. We are hopeful
that algorithmic tricks borrowed from molecular EFM enumeration tools,
such as linear pathway compression, could further decrease memory usage
and run times.

Aside from carbon and nitrogen AEFMs, our work is broadly applicable to
other types of chemical elements. One could feasibly enumerate oxygen or
hydrogen EFMs in small-scale networks, although the significance of these
pathways is unclear for biochemical networks given the ubiquitous nature
of these atoms. We also caution against enumerating hydrogen AEFMs, in
particular, given limitations associated with RXNMapper. Unless explicitly
represented in the reaction SMILES string, RXNMapper does not assign
hydrogen atom mappings. Moreover, RXNMapper cannot map
atoms in reaction SMILES strings longer than 512 characters. Our method
could also be applied to study sulphur metabolism to quantify methionine
or cysteine remodelling, or phosphorus metabolism to quantify nucleotide
synthesis. Outside the realm of biochemistry, our method is applicable to
any periodic table element and possibly relevant for studying general
chemical reaction networks. 

Our AEFM framework opens many new computational problems in metabolism.
While we analyzed one set of fluxes for the HepG2 network, one could
envision a differential AEFM analysis of flux distributions under
differential conditions. These differential AEFMs may reveal changes in
metabolic subsystem activity or predict new pathways of source metabolite
remodelling. For instance, analyzing carbon AEFMs of C1 compounds may be
particularly useful for quantifying or optimizing carbon fixation pathways
in C1-utilizing organisms. Another important question is whether a joint
analysis of AEFM weights could reveal broader metabolic changes within
a condition. We observed that many AEFMs across distinct source
metabolites converged on common internal metabolites. Intuitively, the
weights of AEFMs rooted on different source metabolites, but sharing
common downstream metabolite subsequences, should be related, and these
relationship may further speed AEFM enumeration or weight computations.
Finally, more work is needed to establish relationships between AEFM
decompositions and more traditional EFM decompositions of fluxes, and to
explore the possibility that feasible AEFM calculations might somehow aid
in EFM calculations.

\section*{Materials and Methods}

%\subsection*{Reagents and Tools table}

%See tables/reagents-and-tools-table.xls.

\subsection*{Methods and Protocols}

\subsubsection*{Datasets}

Five highly-curated GEMs were chosen in this study to reflect a diversity
of organisms, metabolites and biochemical reactions. Four networks were
obtained from the BiGG UCSD database~\citep{KingZacharyA2016BMAp} in
ascending order of network size. Datasets were downloaded as SBML (.xml)
files, and the corresponding stoichiometry matrices, metabolites and
reactions were extracted using the Julia SBML.jl package. The HepG2 liver
cancer cell line dataset was chosen for its single set of estimated steady
state fluxes estimated from a FBA with external fluxes constrained by
experimental measurements of amino acid, glucose, pyruvate, and lactate.
These cells were grown in DMEM with 10\% fetal calf serum,
22\SI{}{\milli\molar} glucose, and 1.8\SI{}{\milli\molar}
glutamine~\citep{NilssonAvlant2020Qaoa}. The steady state fluxes were computed
from the FBA model provided by the authors (figure3\_metabolitebalances.m;
MATLAB R2018a) (\citealt{NilssonAvlant2020Qaoa}; Data ref:
\citealt{NilssonAvlant2020QaoaDATASET}). Across all models, we modelled all
reversible reactions as net forward reactions to simplify the number of
molecular and AEFMs as the majority of experimental and computational methods
can only estimate net fluxes. We assumed the reaction stoichiometries of the
BiGG datasets represented the canonical net directions of all reversible
reaction. For the HepG2 dataset, the net directions were chosen based on the
net fluxes estimated from the FBA model.

\subsubsection*{Extracting metabolites SMILES strings}

SMILES structures were not encoded in the BiGG SBML files nor the HepG2
dataset. Thus, we manually extracted SMILES strings corresponding to each
metabolite across all GEMs. Isomeric SMILES strings were taken from
PubChem~\citep{KimSunghwan2023P2u}, the Human Metabolome
Database~\citep{WishartDavidS2022H5tH}, or MetaNetX
database~\citep{MorettiSebastien2021Munf}. Metabolites shared across GEMs
were assigned the same SMILES structure in addition to identical
metabolites stratified across distinct subcellular compartments. Generic
structures such as R-groups were not modelled in this study and were
replaced with a single bonded hydrogen under the assumption this
modification would not alter atom mappings, especially since RXNMapper
does not map hydrogen atoms unless they are explicitly represented in the
SMILES strings. SMILES strings were also not assigned to metabolites with
ambiguous structures nor pseudometabolites with no known chemical
structure. All pseudometabolites and reactions involving pseudometabolites
were removed from the stoichiometry matrix (Table 2).

\subsubsection*{Constructing reaction SMILES strings}

RXNMapper performs atom mapping on a string that encodes the SMILES
structures of all substrates and products for a given reaction. These
reaction SMILES strings are constructed by concatenating SMILES strings of
individual substrates and products. Multiple substrates or products are
separated by the symbol `.', while the reaction arrow delimiting
substrates and products is denoted by the symbol `$>>$'. The substrate and
product SMILES strings are concatenated in order of metabolite appearance
within the stoichiometry matrix, and multiple stoichiometric copies are
consecutively repeated in the reaction SMILES strings. Non-integer
stoichiometric coefficients are therefore not allowed and these
pseudoreactions were removed from the stoichiometry matrix (see Table 2).
To maintain steady state flux in the HepG2 network, flux along
pseudoreactions were replaced with unimolecular source and sink fluxes
consuming each pseudoreaction substrate and producing each pseudoreaction
product. We manually checked several reactions to ensure the assigned atom
mappings were correct. Crucially, we found RXNMapper to fail on nearly all
classes of ATP-independent transaminase reactions. This problem has
previously been reported in an atom mapping benchmark study prior to the
publication of RXNMapper~\citep{PreciatGonzalezGermanA.2017Ceoa}. These
transaminase reactions were identified in all metabolic networks and
manually atom mapped based on the KEGG reaction mappings (Figure \ref{figSte}).

\subsubsection*{Pre-processing genome-scale metabolic models}

Several pre-processing steps were required to compute AEFMs from the GEMs.
The stoichiometric coefficients of external reactions were set to one and
their fluxes rescaled to carry equivalent units of source and sink
metabolites. All reactions containing pseudometabolites with no known
SMILES strings were removed from the stoichiometry matrix and replaced
with unimolecular source and sink reactions carrying equivalent units of
flux. The same procedure was also applied to reactions containing
non-integer stoichiometries. These pseudometabolites were next removed
from the stoichiometry matrix rows, and unimolecular reactions involving
the same external or internal metabolites were aggregated. Finally, we had
to remove a small number of reactions whose reaction SMILES strings
exceeded the 512 token limit of RXNMapper. These reactions were likewise
converted into unimolecular source and sink reactions to maintain steady
state flux across the network (see Table 2). All pre-processing steps
described above are wrapped into convenience functions within our Julia
package MarkovWeightedEFMs.jl.

\subsubsection*{Benchmarking enumeration of atomic-versus-standard EFM}

We benchmarked our Julia package MarkovWeightedEFMs.jl version 2.0
running Julia version 1.10.2 to enumerate
AEFMs versus the molecular EFM enumeration program FluxModeCalculator
(\cite{vanKlinkenJanBert2016Faet}) running MATLAB version R2018a with
default parameters including network compression. We chose
FluxModeCalculator because it is the state-of-the-art program for
enumerating molecular EFMs. Benchmarks were performed on a 64-bit Linux
operating system with 64GB of memory and an AMD 5950X CPU with 16 cores/32
threads. All benchmarks reported the wall time to run the atomic or
molecular EFM enumeration functions and did not include startup times to
open Julia/MATLAB, precompile Julia functions, or load packages and input
data. A benchmark time of `did not finish' (DNF) was assigned if either
program could not complete (A)EFM enumeration within 7 days. In practice,
only FluxModeCalculator failed to enumerate the molecular EFMs in three of
the five datasets.

\subsubsection*{Constructing atomic transition graphs}

For each carbon and nitrogen atom within a given source metabolite, we
identify all possible metabolite/atom states that are traversable by that
initial state and constrained by the reaction atom mapping predictions.
For reactions involving single stoichiometric copies of each substrate and
product, the atom of interest will always move from a single substrate to
product molecule. If this atom is present in a substrate with multiple
stoichiometric copies, we assume there is an equal probability of that
atom occurring in either copy. Depending on the substrate stoichiometric
copy, that atom may move to (i) different product metabolite/atom states,
(ii) the same product but at different atomic positions, or (iii) the same
metabolite/atom position. In all cases, we set the atomic flux of the atom
transitioning from either stoichiometric copy equal to the reaction fluxes
to maintain atomic mass conservation (Figure \ref{figSaccr}). The
resulting atomic transition graph is completed by connecting any sink
fluxes to the root of the tree initialized on the source metabolite/atom
position. 

\subsubsection*{ACHMC model}

Following \citeauthor{ChitpinJustinG2023AMct}, the atomic transition
graphs were analyzed by an atomic version of our CHMC (ACHMC) method to
enumerate the corresponding AEFMs and compute their weights. In this
formulation, each ACHMC state corresponds to a sequence of metabolite-atom
positions transited by a given source metabolite atom. Each simple cycled
identified by the ACHMC corresponds to an AEFM rooted on a user-specified
source metabolite atom. For the HepG2 flux network, we computed the
probabilities of each AEFM by steady state analysis of the ACHMC. These
AEFM probabilities were scaled to weights such that the totality of
source-to-sink AEFMs was equal to the unimolecular input flux producing
that source metabolite.

\section*{Data availability}

Our software package MarkovWeightedEFMs.jl is available at
\url{https://github.com/jchitpin/MarkovWeightedEFMs.jl}. All source codes
to reproduce this study are available at
\url{https://www.github.com/jchitpin/reproduce-efm-paper-2024}.

\section*{Acknowledgements}

This research was supported by the Natural Sciences and Engineering
Research Council of Canada (NSERC), Discovery grant RGPIN/06604-2019 and
a Digital Research Alliance of Canada (DRAC) Resources for Research Groups
grant to TJP (RRG-TPERKINS). J.G.C. was supported by an NSERC CREATE
Matrix Metabolomics Scholarship, NSERC Alexander Graham Bell Canada
Graduate Scholarship, and Ontario Graduate Scholarship.

\section*{Author contributions}

Justin G. Chitpin: Conceptualization; data curation; methodology;
software; formal analysis; investigation; visualization; writing
- original draft; writing - review and editing.\\

\noindent Theodore J. Perkins: Conceptualization; methodology;
supervision; funding acquisition; investigation; project administration;
writing - review and editing.

\section*{Disclosure and competing interests statement}

The authors declare no competing interests.

\bibliography{biblio-aefms}


\newpage
\pagenumbering{gobble} 
%\section{Figures}
 
%\noindent 
\begin{figure}
\includegraphics[width=\textwidth]{figure1-pipeline.pdf}
\caption{Atomic cycle-history Markov chain (ACHMC) pipeline. (a) An
example molecular EFM in glycolysis. (b) An example AEFM tracing the flow
of a given glucose carbon in glycolysis. (c) The major steps within our
pipeline to enumerate AEFMs and compute their weights. \label{fig1Pipe}}
\end{figure}

\begin{figure}
\includegraphics[width=\textwidth]{figure2-benchmarking.pdf}
\caption{AEFM enumeration is computationally tractable
compared to molecular EFMs in five GEMs. (a) Number of atomic and
molecular EFMs in the five GEMs. (b) The run time of our
MarkovWeightedEFMs.jl method versus FluxModeCalculator running in serial.
(c) Serial run time scaling of MarkovWeightedEFMs.jl as a function of the
number of AEFMs. (d) Speedup of selected ACHMCs as a function of
additional threads. \label{fig2Bench}}
\end{figure}

\begin{figure}
\centerline{\includegraphics[height=8in]{figure3-structural.pdf}}
\caption{Structural analysis of AEFMs in five genome-scale
metabolic models. (a-b) Classes of carbon and nitrogen EFMs. (c-d)
Fraction of carbon and nitrogen AEFM lengths. (e-f) Ratio of the greatest
and smallest number of AEFMs between carbons and nitrogens within the same
source metabolite. Representative metabolites from each dataset are
labelled; labels with asterisks indicate that metabolite is also present
in other networks. \label{fig3Struct}}
\end{figure}


\begin{figure}
\includegraphics[width=\textwidth]{figure4-hepg2-weights.pdf}
\caption{Carbon AEFM weight analysis of the HepG2 metabolic
flux network. (a) A small number of highly active AEFMs explain the
majority of glutamine flow. (b) Subnetwork of the top five AEFMs across
each glutamine carbon collapsed onto metabolites. (c-e) Sankey diagrams of
the top five glutamine carbon AEFMs. (f) Schematic mapping AEFMs onto
traditionally defined metabolic pathways/subsystems. All metabolite names are
labelled following the BiGG database nomenclature. \label{fig4Hepg2}}
\end{figure}

%\section*{Tables and their legends}

% Table 1
\begin{table}
\caption{Summary statistics for the five metabolic networks in this study.}
~\\
\begin{tabular}{l|ccccc}
& \multicolumn{5}{c}{{Carbon metabolism}} \\
{GEM (\# mets $|$ \# rxns)} & {Sources} & {Sinks} & {States} & {Transitions} & {AEFMs} \\ \hline
E. coli core (71 $|$ 144) & 184 & 255 & 14659 & 24355 & 6519 \\
iAB\_RBC\_283	 (274 $|$ 465) & 954 & 1166 & 37122 & 51220 & 7958 \\
iIT341 (433 $|$ 759) & 1486 & 2134 & 86342 & 120545 & 30633 \\
iSB619 (547 $|$ 974) & 2188 & 2986 & 169877 & 246712 & 59444 \\
HepG2 (435 $|$ 545) & 376 & 939 & 9092000 & 11615891 & 1377937\\ 
~\\
& \multicolumn{5}{c}{{Nitrogen metabolism}} \\
{GEM (\# mets $|$ \# rxns)} & {Sources} & {Sinks} & {States} & {Transitions} & {AEFMs} \\ \hline
E. coli core (71 $|$ 144) & 30 & 40 & 246 & 448 & 204 \\
iAB\_RBC\_283	 (274 $|$ 465) & 161 & 214 & 6409 & 9275 & 2091\\
iIT341 (433 $|$ 759) & 314 & 447 & 27440 & 39761 & 10440 \\
iSB619 (547 $|$ 974) & 353 & 516 & 36661 & 54165 & 14214 \\
HepG2 (435 $|$ 545) & 72 & 190 & 19137 & 27860 & 5105 \\
\end{tabular}
\end{table}

%Table 2:
\begin{table}
\caption{Number of metabolites and reactions removed from each network during pre-processing.}
\begin{tabular}{l|ccccc}
 & {E. coli core} & {iAB\_RBC\_283} & {iIT341} & {iSB619} & {HepG2} \\ \hline
 ~\vspace{-1ex}\\
{\small Rxns with non-integer stoichiometries} & 2 & 1 & 3 & 34 & 8 \\
{\small Pseudometabolites with no known structures} & 1 & 68 & 52 & 108 & 33 \\
{\small Rxns with pseudometabolites} & 2 & 142 & 93 & 204 & 46 \\
{\small Rxns exceeding RXNMapper token limit} & 0 & 1 & 4 & 14 & 3 
\end{tabular}
\end{table}

%\section*{Expanded View Figure legends}
\setcounter{figure}{0} 
\renewcommand{\thefigure}{S\arabic{figure}}

\begin{figure}
\centerline{\includegraphics[height=6.25in]{figureS5-atomic-chmc-cartoon-rules.pdf}}
\caption{Rules for assigning atomic fluxes from
atom-mapped reactions with biological examples. Schematic node colours
represent ground truth atom mappings across substrate(s) and product(s).
Node labels correspond to distinct atomic indices of substrate and product
stoichiometric copies.
(a) For reactions with single stoichiometric copies of substrates and
products, each atom within a given substrate moves to a single atomic
position within a product with atomic flux equal to the reaction flux.
(b-d) For reactions with $s$ stoichiometric copies of a given substrate,
there are $s$ copies of each metabolite-atom state on the substrate and
product side of the reaction. The probability of any atom within
a substrate stoichiometric copy moving to a product metabolite-atom state
remains equal to the reaction flux. (b) The stoichiometric copies of the
substrate fuse together into a single product. (c) The stoichiometric
copies of the substrate form distinct products. (d) The $s$ stoichiometric
copies of the substrate form $s$ copies of products with identical atomic
backbones. (e) Biological example of rule a where the atoms in each
substrate map to a single position on the corresponding product. (f)
Biological example of rule b where two stoichiometric copies of
glutathione fuse together to form a single stoichiometric copy of
glutathione disulfide. (g) Biological example of rule c where two
stoichiometric copies of acetyl-CoA undergo a reaction to produce
a stoichiometric copy of two distinct products. \label{figSaccr}}
\end{figure}

%Expanded View 1
\begin{figure}
\centerline{\includegraphics[height=7.5in]{figureS1-gem-hairballs.pdf}}
\caption{Network visualization of the five GEMs after
pre-processing. Metabolites (non-grey nodes) are connected to reaction
entities (grey nodes) with node sizes proportional to their degree. Graphs
were constructed in Gephi version 0.10 using the Yifan-Hu proportional
network layout with default parameters. (a) E. coli core. (b) iAB RBC 283.
(c) iIT341. (d) iSB619. (e) HepG2. \label{figSgh}}
\end{figure}

%Expanded View 2
\begin{figure}
\centerline{\includegraphics[width=3.5in]{figureS2-benchmarking-fluxmodecalculator-parallel.pdf}}
\caption{FluxModeCalculator scaling across multiple threads. \label{figSbfp}}
\end{figure}

%Expanded View 3
\begin{figure}
\centerline{\includegraphics[width=\textwidth]{figureS3-atomic-state-space.pdf}}
\caption{Most ACHMCs traverse less than 10\% of the
total metabolite-carbon/nitrogen state space. Each point corresponds to an
ACHMC rooted on a given source metabolite carbon or nitrogen. \label{figSass}}
\end{figure}

\begin{figure}
\centerline{\includegraphics[width=\textwidth]{figureS4-cumulative-combined-legend.pdf}}
\caption{Cumulative explained carbon mass flow of input
glucose and amino acid carbons from the HepG2 flux network. The cumulativeexplained mass flow is truncated at 99\% to avoid displaying AEFMs that carry little atomic mass flow. \label{figSccl}}
\end{figure}



\begin{figure}
\centerline{\includegraphics[height=8in]{figureS6-transaminase-examples.pdf}}
\caption{Examples of incorrectly atom-mapped transaminase
reactions present in one or more of the five networks. (a) Phenylalanine
transaminase. (b) Isoleucine transaminase. (c) Tyrosine transaminase. (d)
Aspartate transaminase.\label{figSte}}
\end{figure}


\end{document}

